\documentclass[../main]{subfiles}
%\usepackage[spanish]{babel}
\usepackage[lining,light]{InriaSans}
\usepackage{tcolorbox}
\usepackage{qrcode}
\usepackage{hyperref}
\usepackage{tikz}
\usepackage{geometry}
\usepackage[pages=all]{background}
\usepackage{svg}
\usepackage{relsize}
\tikzset{fontscale/.style = {font=\relsize{#1}}}
\renewcommand*\oldstylenums[1]{{\fontfamily{InriaSans-OsF}\selectfont #1}}
\let\oldnormalfont\normalfont
\def\normalfont{\oldnormalfont\mdseries}

\renewcommand{\baselinestretch}{1.5}

\geometry{
 a4paper,
 left=30mm,
 top=15mm,
 right=15mm,
 bottom=55mm
 }
 
\makeatletter
\newcommand*{\currentname}{\@currentlabelname}
\makeatother
\backgroundsetup{scale=1, color=black, opacity=1, angle=0, 
 contents={
\begin{tikzpicture}[remember picture, overlay, shift={(-10.5,-14.85)}]
 \tikzset{every node}=[font=\small\sffamily]
\node at (4.3,2.75){\includegraphics[width=3.5cm]{\logo}};
%\draw [step=0.5,gray, thin] (2.5,4.5) grid (20,28.5) ;
\draw [line width=0.6mm,black] (2.5,1) -- (20,1) -- (20,28.7) -- (2.5,28.7) -- (2.5,1);
\draw [line width=0.6mm,black] (2.5,4.5) -- (20,4.5);
\draw [line width=0.6mm,black] (6,1) -- (6,4.5);
\draw [line width=0.6mm,black] (20,2.75) -- (6,2.75);
\draw [line width=0.6mm,black] (13,2.75) -- (13,4.5);
\draw [line width=0.6mm,black] (16.5,2.75+2*0.58) -- (16.5,1);
\draw [line width=0.6mm,black] (13,2.75+0.58) -- (20,2.75+0.58);
\draw [line width=0.6mm,black] (13,2.75+2*0.58) -- (20,2.75+2*0.58);
\draw [line width=0.6mm,black] (16.5,1.58) -- (20,1.58);
\draw [line width=0.1mm,black] (2.5,14.85) -- (2.3,14.85);
\node at (17.6,1.3){HOJA N.$^o$:\thepage};
\node at (17.2,2.4){V.$^o$B.$^o$:};
\node at (14,4.2){NOMBRE:};
\node at (13.8,3.6){CURSO:};
\node at (14.8,3){FECHA: .../.../......};
\node at (17.3,3.6){T.P.N.$^o$:\TP};
\node at (17.15,3){ESC.:};
\node at (9.5,3.7) {\materia};
\node at (12.1,1.9) {\currentname};
\node at (6.9,1.87) {\qrcode[height=1.5cm]{\urlqr}};
\end{tikzpicture}
 }
}

\begin{document}
\pagestyle{empty}
\info{PROGRAMA}{
CICLO LECTIVO 2025

ÁREA: TÉCNICA

ESPACIO CURRICULAR: Dibujo Técnico 			

CURSO: 2

PROFESOR/A: Ferreyra Gustavo 
}
\setcounter{chapter}{1}
\section{Programa}
\subsection{Bloque 1}
\begin{itemize}
    \item Letras formato de láminas.
    \item Líneas, aplicaciones.
    \item Interrupciones de cuerpos.
    \item Líneas de rotura, aplicaciones.
    \item Rayados y colores convencionales.
    \item Modo de plegar los planos.
    \item Rotulación de planos.
\end{itemize}
\subsection{Procedimientos 1}
\begin{itemize}
    \item Representar un cuerpo aplicando distintos tipos de líneas.
    \item Presentación de dibujos en láminas con formato y rótudoadecuado.
    \item Plegado de planos.
\end{itemize}
\subsection{Bloque 2}
\begin{itemize}
    \item Métodos de representación: sistema IRAM.
    \item Sistemas de representación: perspectiva caballera, perspectiva isométrica.
    \item Rebatimiento y desarrollos.
\end{itemize}
\subsection{Procedimientos 2}
\begin{itemize}
    \item Representar un cuerpo por medio de vistas, en los distintos métodos de representación.
    \item Dadas las vistas principales de un cuerpo, identificar en que método está representado.
    \item Dibujar un cuerpo en perspectiva dadas sus vistas principales.
\end{itemize}
\subsection{ Bloque 3}
\begin{itemize}
    \item Acotaciones.
    \item Acotaciones de los dibujos.
    \item Ejecución de las líneas de cota y de referencia.
    \item Sistema de acotamiento.
    \item Acotaciones condicionadas a lineas base o de referencia.
\end{itemize}
\subsection{Procedimientos 3}
\begin{itemize}
    \item Acotado de un dibujo.
    \item Presentación de dibujos en láminas acotados y en escala correcta.
    \item confección de croquis.
\end{itemize}
\subsection{Bloque 4}
\begin{itemize}
    \item Generalidades de software FreeCad y AutoCAD.
    \item Estructura de la pantalla.
    \item Símbolo de vista.
    \item Modos de seleccionar objetos.
    \item Modos de captura de objetos.
    \item comandos: Línea, rectángulos, polígonos, copiar, pegar, recortar, empalmar líneas, borrar, extender una línea, escribir un texto, rotar, espejar, trazar paralelas.
    \item Editar texto.
\end{itemize}
\subsection{Procedimientos 4}
\begin{itemize}
                \item Dibujo de formato de hoja y rótulo, en software freeCAD.
                \item Ejecución en FreeCAD, de los dibujos de la guía de trabajos prácticos.
                \item Impresión de planos realizados.
                \item Aplicación de conocimientos en el resto de los bloques.
                \item Uso de plotter e impresora.
\end{itemize}
\end{document}
